\documentclass[10pt,landscape,letterpaper]{article}
\usepackage[landscape,left=3mm,right=3mm,top=3mm,bottom=3mm]{geometry}
\usepackage{multicol}
\usepackage{amsmath,amssymb,amsfonts}
\usepackage{enumitem}
\usepackage{xcolor}
\usepackage{titlesec}
\usepackage{fontspec}
\usepackage{booktabs}
\usepackage{graphicx}

% Font settings
\setmainfont{Verdana}
\setsansfont{Verdana}

% Colors
\definecolor{sectionblue}{RGB}{0,51,102}
\definecolor{conceptcyan}{RGB}{0,139,139}
\definecolor{processpurple}{RGB}{128,0,128}
\definecolor{categorygreen}{RGB}{0,128,0}
\definecolor{highlightyellow}{RGB}{255,255,150}

% Font size and spacing
\renewcommand{\normalsize}{\fontsize{5pt}{5pt}\selectfont}
\normalsize

% Section formatting
\titleformat{\section}{\color{sectionblue}\bfseries\fontsize{6pt}{6pt}\selectfont}{}{0em}{}
\titleformat{\subsection}{\color{conceptcyan}\bfseries\fontsize{5.5pt}{5.5pt}\selectfont}{}{0em}{}
\titlespacing*{\section}{0pt}{1pt}{0.5pt}
\titlespacing*{\subsection}{0pt}{0.5pt}{0.25pt}

% List settings
\setlist[itemize]{leftmargin=*,nosep,topsep=0pt,parsep=0pt,partopsep=0pt,itemsep=0pt}
\setlist[enumerate]{leftmargin=*,nosep,topsep=0pt,parsep=0pt,partopsep=0pt,itemsep=0pt}

% Paragraph settings
\setlength{\parindent}{0pt}
\setlength{\parskip}{0.3pt}
\setlength{\columnsep}{3mm}

\pagestyle{empty}

\newcommand{\concept}[1]{\textcolor{conceptcyan}{\textbf{#1}}}
\newcommand{\process}[1]{\textcolor{processpurple}{\textbf{#1}}}
\newcommand{\category}[1]{\textcolor{categorygreen}{\textbf{#1}}}
\newcommand{\important}[1]{\colorbox{highlightyellow}{#1}}

\begin{document}
\begin{multicols}{6}

\section{SCALABILITY}
\concept{Scalable}: system adapts to increased/reduced demand w/o redesign.
\category{Horizontal}: add machines.
\category{Vertical}: add resources to one machine.\\
Scalability limited by \important{bottlenecks}.\\
$4\times$ nodes $\approx$ $4\times$ capacity (efficient growth).

\subsection{Speedup}
$\text{Speedup} = \frac{\text{time w/ 1 core}}{\text{time w/ N cores}}$\\
Ideal: linear. Reality: sublinear (serial parts + overhead).

\subsection{Amdahl's Law}
\important{$S(N) = \frac{1}{(1-f) + f/S_{\text{part}}}$}\\
$f$ = parallelizable fraction\\
Even $\infty$ cores can't beat $\frac{1}{1-f}$\\
More parallelism can \textcolor{red}{decrease} perf beyond sweet spot.

\subsection{Granularity}
\category{Coarse-grain}: large tasks/core (efficient)\\
\category{Fine-grain}: small tasks/core (too much overhead)

\subsection{Dependencies}
Tasks may depend on others $\to$ limits parallelism.\\
Max speedup = \concept{critical path} (longest chain).\\
\concept{Embarrassingly parallel}: no dependencies.

\section{RACE CONDITIONS}
\concept{Race condition}: result depends on exact timing of concurrent threads.\\
Cause: concurrent \textbf{writes} to shared state.\\
Read-only access is safe.

\subsection{Consistency Models}
\category{Sequential}: result = valid sequential interleaving; each core's ops in program order\\
\category{Strong}: after write, all reads return new value\\
\category{Weak}: no guarantee until condition met\\
\category{Eventual}: if no more writes, reads \emph{eventually} return latest\\
Variants: causal, read-your-writes, monotonic

\section{MUTUAL EXCLUSION}
\concept{Critical section}: code where concurrent access causes problems.\\
\concept{Mutual exclusion}: only one thread in critical section at a time. Reduces parallelism.

\subsection{Locking}
\process{LOCK(X)}: if unlocked, lock \& continue; if locked, wait.\\
\process{UNLOCK(X)}: release lock.\\
Granularity: one big lock (simple) vs.\ per-variable (more parallel).

\subsection{Deadlocks}
Two+ threads each hold lock other needs $\to$ neither proceeds.\\
Detect via \concept{wait-for graph}: cycle = deadlock.\\
\category{Solutions}:
\begin{itemize}
\item Lock manager (centralized, limits scalability)
\item Resource hierarchy (fixed global order)
\item Chandy/Misra (dirty/clean fork protocol)
\end{itemize}
\concept{Optimistic concurrency}: allow concurrent updates, merge; if conflict, serialize.

\section{MEMORY ARCHITECTURES}
\subsection{SMP}
\concept{Symmetric Multiprocessing}: all cores share memory, uniform latency.\\
Simple, easy load balance. Limited $\sim$10 cores.

\subsection{NUMA}
\concept{Non-Uniform Memory}: memory local to each CPU.\\
Any core accesses any byte, but local 2--3$\times$ faster.\\
Better scalability, harder programming.

\subsection{Shared-Nothing}
Independent machines on network.\\
Each CPU: local memory only; remote = message passing.\\
Best scalability, nontrivial programming.

\section{SCALING DISTRIBUTED}
\subsection{Load Balancing}
Distribute TCP flows (not packets!) across frontends.\\
\process{Layer-4 LB}: hash 5-tuple $\to$ machine.

\subsection{Frontend / Backend}
Separate request handling (FE) from data (BE).\\
FE stateless $\to$ easy to scale.\\
Problem: replicated FE + local data = no consistency.\\
Solution: centralize data in BE.

\subsection{Sharding}
Split BE data across machines.\\
E.g.\ accounts 0--4 on server A, 5--9 on B.\\
FE routes to correct BE. Scales read \& write.

\subsection{Cost of Scalability}
\important{Scalability rarely free.}\\
Single-thread SSD beat Spark 128-core for PageRank.\\
(McSherry et al.\ ``Scalability! But at what COST?'')

\section{THE WEB}
\concept{WWW}: collection of interconnected docs running on top of the Internet (not = Internet).\\
\concept{Ingredients}: HTML, URIs/URLs, DNS, HTTP, client/server, browser, scripting, cookies.

\subsection{URIs / URLs}
\concept{URI}: umbrella (URN + URL)\\
\concept{URN}: identifies \emph{what} (\texttt{urn:isbn:059652})\\
\concept{URL}: identifies \emph{where}\\
\important{\texttt{<scheme>://[user@]<host>[:port]/[path][?params]}}

\section{CLIENT-SERVER MODEL}
\concept{Servers}: offer services.\\
\concept{Clients}: consume via request/response.\\
\concept{Peer-to-peer}: no special machines, all peers equal.\\
Well-known ports: 80 (HTTP), 443 (HTTPS), 22 (SSH).\\
\process{Server challenges}: concurrency, fault tolerance, state, consistency, security.

\section{HTTP PROTOCOL}
\concept{HTTP}: text-based request/response on TCP.\\
Port 80 (HTTP), 443 (HTTPS w/ TLS). First spec 1990.

\subsection{Request}
\texttt{METHOD URI HTTP/1.1}\\
+ headers + blank line + body (optional)

\subsection{Methods}
\begin{tabular}{@{}ll@{}}
\concept{GET} & Retrieve resource \\
\concept{HEAD} & Metadata only \\
\concept{PUT} & Store at URI \\
\concept{DELETE} & Remove resource \\
\concept{POST} & Submit data (forms)
\end{tabular}

\subsection{Headers}
\concept{Connection}: close / keep-alive\\
\concept{Content-Length}: bytes in body\\
\concept{Content-Type}: MIME type

\subsection{Status Codes}
\begin{tabular}{@{}ll@{}}
\important{200} & OK \\
\important{301} & Moved Permanently \\
\important{304} & Not Modified \\
\important{401} & Unauthorized \\
\important{403} & Forbidden \\
\important{404} & Not Found \\
\important{500} & Internal Server Error
\end{tabular}

\subsection{HTTPS}
HTTP + TLS. Encrypts traffic. Client verifies server via cert. IP still visible.

\section{HTTP/2 \& HTTP/3}
\subsection{HTTP/1.1 Problems}
\begin{itemize}
\item \concept{Head-of-line blocking}: slow object blocks later ones
\item No prioritization
\item No server push
\end{itemize}

\subsection{HTTP/2 (2015)}
\category{Binary framing}: msgs as binary frames (not text)\\
\category{Multiplexing}: single TCP conn, multiple streams\\
\category{Priority}: weights + parent deps\\
\category{HPACK}: header compression\\
\category{Server push}: PUSH\_PROMISE; client can RST\_STREAM

\subsection{HTTP/3}
Replaces TCP w/ \concept{QUIC} (RFC 9000, on UDP).\\
Streams fully independent. Loss in one $\neq$ block others.

\section{HTML \& CSS}
\concept{HTML}: structure + semantic content.\\
Elements: \texttt{<h1>}-\texttt{<h6>}, \texttt{<p>}, \texttt{<br>}, \texttt{<ul>/<li>}, \texttt{<img>}, \texttt{<a>}, \texttt{<table>}.\\
\concept{CSS}: separate content from formatting.\\
\texttt{<div>}/\texttt{<span>}: generic containers.\\
\concept{Cascade}: multiple sheets, defined precedence.

\subsection{Forms}
\texttt{<form action="url" method="get|post">}\\
\category{GET}: params in URL, idempotent, length-limited\\
\category{POST}: params in body, binary OK, browser warns resend

\section{DYNAMIC CONTENT}
\subsection{CGI}
\concept{CGI}: server runs external program per request.\\
Program reads env vars + stdin, writes HTML stdout.\\
Drawbacks: new \textbf{process}/req (slow), security risks.

\subsection{Java Servlets}
Java class in servlet container (Tomcat).\\
One instance; each req = own \textbf{thread} (not process).\\
\begin{tabular}{@{}ll@{}}
 & \concept{CGI} / \concept{Servlets} \\
Handled by & processes / threads \\
State & filesystem / container \\
Security & risky / Java sandbox
\end{tabular}

\subsection{Routing}
Map URL patterns to handlers.\\
\concept{Spark}: \texttt{get("/hello", (req,res) -> "Hi")}\\
Patterns: wildcards, path params, query params.

\section{SESSIONS \& COOKIES}
\concept{Session}: sequence of related interactions.\\
State kept server-side; client gets \concept{session ID}.\\
\important{Session ID must be large + random!}

\subsection{Cookies}
Small KV pairs in browser ($\leq$4KB).\\
Server: \texttt{Set-Cookie} header.\\
Browser sends back on matching domain+path.\\
Attrs: expiration, domain, path.

\section{TRACKING \& ADS}
\concept{Ad networks}: combine tracking across sites.\\
\concept{Third-party cookies}: embedded content carries ad domain's cookies $\to$ cross-site tracking.\\
Infer: gender, age, interests, income.\\
\important{Google Ads: \$147B revenue (2020).}\\
Privacy: browsers blocking 3rd-party cookies.\\
Alternatives: micropayments, FLoC/FLEDGE (cohorts).

\end{multicols}
\end{document}
